\section{근, 중복도와 형식적 미분}
\label{sec:roots-derivatives}

\begin{learninggoal}
근의 개수와 차수의 관계를 이해하고, 형식적 미분을 이용해 중근을 판정한다.
\end{learninggoal}

\begin{definition}
체 $K$와 영이 아닌 다항식 $f\in K[X]$에 대해 $\alpha\in K$가 $f(\alpha)=0$을 만족하면 $\alpha$를 $f$의 근이라 한다. $(X-\alpha)^m\mid f$이고 $(X-\alpha)^{m+1}\nmid f$인 유일한 양의 정수 $m$을 $\alpha$의 중복도라 한다.
\end{definition}

영다항식은 모든 $\alpha$에서 값이 $0$이지만 각 선형인수의 모든 거듭제곱으로 나누어지므로 유한한 중복도를 정의할 수 없다. 따라서 근의 중복도를 다룰 때에는 항상 다항식이 영이 아니라고 가정한다.

\begin{theorem}[근의 개수 정리]
영이 아닌 $f\in K[X]$의 차수가 $n$이면 $K$에 속하는 서로 다른 근들의 중복도 합은 많아야 $n$이다.
\end{theorem}

\begin{proof}
근 $\alpha$가 하나 있으면 인수정리에 의해 $f=(X-\alpha)g$이고 $\deg g=n-1$이다. 선형인수 하나를 제거할 때마다 해당 근의 중복도와 전체 차수가 각각 $1$씩 감소한다. 이를 차수에 대해 귀납하면 서로 다른 근들의 중복도 합이 $n$을 넘을 수 없다.
\end{proof}

\begin{corollary}
$K$가 무한체이고 두 다항식 $f,g\in K[X]$가 모든 $a\in K$에서 $f(a)=g(a)$를 만족하면 $f=g$이다.
\end{corollary}

\begin{proof}
$f-g$가 모든 $a\in K$를 근으로 갖는다. 영이 아닌 다항식은 차수보다 많은 서로 다른 근을 가질 수 없으므로 $f-g=0$이다.
\end{proof}

\begin{definition}
다항식 $f=\sum_{i=0}^n a_iX^i$의 \emph{형식적 도함수}를
\[
 f'=\sum_{i=1}^n ia_iX^{i-1}
\]
로 정의한다. 이는 모든 표수에서 정의되며
\[
 (f+g)'=f'+g',\qquad (fg)'=f'g+fg'
\]
를 만족한다.
\end{definition}

\begin{theorem}[중근 판정]
영이 아닌 $f\in K[X]$이고 $\alpha\in K$가 $f$의 근이라 하자. 다음은 동치이다.
\begin{enumerate}[label=(\roman*)]
  \item $\alpha$의 중복도는 적어도 $2$이다.
  \item $f'(\alpha)=0$이다.
  \item $(X-\alpha)$가 $f$와 $f'$를 모두 나눈다.
\end{enumerate}
\end{theorem}

\begin{proof}
$f=(X-\alpha)^mg$와 $g(\alpha)\ne0$을 만족하는 중복도 $m\ge1$을 택하면
\[
 f'=m(X-\alpha)^{m-1}g+(X-\alpha)^mg'.
\]
$m=1$이면 $f'(\alpha)=g(\alpha)\ne0$이고, $m\ge2$이면 두 항 모두 $X-\alpha$로 나누어져 $f'(\alpha)=0$이다. 따라서 (i)와 (ii)는 동치이다. 인수정리를 $f'$에 적용하면 (ii)와 (iii)의 동치도 얻는다.
\end{proof}

\begin{corollary}
영이 아닌 $f\in K[X]$가 어떤 확대체에서도 중근을 갖지 않을 필요충분조건은
\[
  \gcd(f,f')=1
\]
이다. 이런 다항식을 \emph{제곱인수가 없는 다항식}이라 한다.
\end{corollary}

\begin{proof}
$d=\gcd(f,f')$라 하자. $\deg d>0$이면 $d$는 어떤 분해체에서 근 $\alpha$를 갖는다. 그러면 $f(\alpha)=f'(\alpha)=0$이므로 중근 판정에 의해 $\alpha$는 $f$의 중근이다.

반대로 어떤 확대체 $L/K$에서 $\alpha$가 $f$의 중근이라고 하자. 그러면 $f(\alpha)=f'(\alpha)=0$이다. 만약 $\gcd(f,f')=1$이면 베주 항등식으로
\[
  af+bf'=1
\]
인 $a,b\in K[X]$가 존재한다. 이 식을 $L$에서 $\alpha$에 대입하면 $0=1$이라는 모순을 얻는다. 따라서 $\gcd(f,f')\ne1$이다.
\end{proof}

\begin{example}
$\Q[X]$에서 $f=X^3-3X+2=(X-1)^2(X+2)$이다. 도함수는
\[
 f'=3X^2-3=3(X-1)(X+1)
\]
이므로 $\gcd(f,f')=X-1$이고 $1$이 중근임을 확인할 수 있다.
\end{example}

\subsection{양의 표수에서의 주의점}

표수가 $p>0$이면 $p\cdot1_K=0$이므로
\[
 (X^p)'=pX^{p-1}=0.
\]
따라서 비상수 다항식의 도함수가 $0$일 수 있다.

\begin{proposition}
$K$의 표수가 $p>0$일 때 $f'=0$일 필요충분조건은 어떤 $g\in K[X]$에 대해
\[
 f(X)=g(X^p)
\]
로 쓸 수 있는 것이다.
\end{proposition}

\begin{proof}
$f=\sum_i a_iX^i$이고 $f'=0$이라고 하자. $p\nmid i$이면 $i\cdot1_K$는 영이 아닌 체의 원소이므로 단위원이다. 따라서
\[
  ia_i=0\quad\Longrightarrow\quad a_i=0.
\]
즉 $p$의 배수인 지수만 남으므로
\[
  f(X)=\sum_j a_{pj}X^{pj}=g(X^p),
  \qquad g(T)=\sum_j a_{pj}T^j
\]
로 쓸 수 있다. 반대로 $f(X)=g(X^p)$이면 모든 비영 항의 지수가 $p$의 배수이므로 형식적 도함수는 $0$이다.
\end{proof}

\begin{pitfall}
실해석의 미분과 달리 형식적 미분에는 극한이 없다. 다만 계수와 지수에 대한 대수적 규칙만 사용한다. 특히 양의 표수에서는 도함수가 $0$이라고 해서 다항식이 상수인 것은 아니다.
\end{pitfall}
